\documentclass[11pt]{article}
\usepackage[margin=0.82in]{geometry}
\usepackage{booktabs}
\usepackage{hyperref}

\title{HAOS-IIP Ferron Extension: Testable Signatures}
\author{Tomislav Rupi\'c \\ Independent Researcher}
\date{May 2026}

\begin{document}
\maketitle

\section*{Scope}

This note converts Materials Bridge Line A into forward-facing, falsifiable
prediction records. It is not a new theory, not a proof of HAOS-IIP, and not a
claim that ferrons embody HAOS-IIP.

\section*{Calibration Source}

Release \texttt{64.1} provides the bounded calibration point: real Figshare
ferron / NbOI$_2$ data loads; \texttt{15/15} target-window spectral records are
found near \texttt{3.13 THz}; mean peak frequency is \texttt{3.13893 THz}; mean
spectral recoverability is \texttt{0.928455}; published Fig. 2d target-band
traces contain \texttt{8} traces and \texttt{216} time points; the published
group-velocity proxies are \texttt{58.8 km/s} at \texttt{96 nm} and
\texttt{100 km/s} at \texttt{224 nm}; the computed-from-time-trace STFT
diagnostic gives mean envelope correlation \texttt{0.995318}; and the unified
persistence proxy is \texttt{0.957192}.

The raw-grid boundary remains active:
\texttt{raw\_stft\_time\_frequency\_grid\_status=NO\_STFT\_DATA\_FOUND}.
Published target-band traces and computed STFT diagnostics are not raw
time-frequency maps.

\section*{Field-Facing Predictions}

\begin{center}
\begin{tabular}{p{0.25\linewidth}p{0.35\linewidth}p{0.30\linewidth}}
\toprule
Prediction & Testable readout & Falsifier \\
\midrule
Polar-mode address selection &
Layered polar ferroelectrics or related polar oxyhalides should show a
narrowband coherent readout in a pre-registered target window around an
independently measured polar mode. &
If matched polar materials show only broadband emission, or if non-polar
controls show the same target-window recoverability, the prediction fails for
that material class. \\
\addlinespace
Thickness / boundary propagation dependence &
Target-band packet arrival time should increase with pump-probe separation, and
group velocity should vary with thickness or confinement under matched
excitation. &
If controlled thickness series show no distance-ordered peak delay and no
thickness dependence beyond uncertainty, the prediction fails for that material
class. \\
\addlinespace
Re-addressing after secondary perturbation &
After a secondary pulse, mild defect scattering, or boundary perturbation below
damage threshold, the readout should preferentially return to the original
target frequency window. &
If secondary perturbation repeatedly produces broadband thermalization or a new
dominant frequency with no preferential recovery of the original target window,
the prediction fails for that perturbation regime. \\
\addlinespace
Cross-dataset narrowband control discrimination &
Pre-registered target-window recoverability should separate coherent
directional propagation records from broadband or non-propagating controls
better than amplitude-only ranking. &
If target-window recoverability does not distinguish coherent propagation
records from controls better than amplitude-only baselines, the prediction
fails for the selected dataset class. \\
\bottomrule
\end{tabular}
\end{center}

\section*{Provisional Thresholds}

Suggested first-pass thresholds are intentionally conservative and revisable:
target peak in at least \texttt{80\%} of controlled records, linewidth proxy
below \texttt{0.1 THz} where estimable, mean spectral recoverability at or above
\texttt{0.85}, frequency stability at or above \texttt{0.95}, and target-band
recoverability remaining above \texttt{0.70} before visible propagation failure.

\section*{Boundary}

These are field-facing diagnostic candidates. They do not claim new physics, a
materials ontology, a universal harmonic grid, or proof of HAOS-IIP. The
numbers above are provisional thresholds seeded by the ferron audit, not
universal constants.

\end{document}
